\documentclass[10pt,a4paper]{amsart}
\usepackage{amsmath,amssymb,amsthm,graphicx}
\usepackage{verbatim}
\newenvironment{reference}{\comment}{\endcomment}
%\newenvironment{reference}{}{}
\newenvironment{slogan}{\comment}{\endcomment}
\newenvironment{history}{\comment}{\endcomment}

% For commutative diagrams we use Xy-pic
\usepackage[all]{xy}

% We use 2cell for 2-commutative diagrams.
\xyoption{2cell}
% Theorem environments.
%
\theoremstyle{plain}
\newtheorem{theorem}[subsection]{Theorem}
\newtheorem{proposition}[subsection]{Proposition}
\newtheorem{lemma}[subsection]{Lemma}

\theoremstyle{definition}
\newtheorem{definition}[subsection]{Definition}
\newtheorem{example}[subsection]{Example}
\newtheorem{exercise}[subsection]{Exercise}
\newtheorem{situation}[subsection]{Situation}

\theoremstyle{remark}
\newtheorem{remark}[subsection]{Remark}
\newtheorem{remarks}[subsection]{Remarks}

\numberwithin{equation}{subsection}

% Macros
%
\def\lim{\mathop{\mathrm{lim}}\nolimits}
\def\colim{\mathop{\mathrm{colim}}\nolimits}
\def\Spec{\mathop{\mathrm{Spec}}}
\def\Hom{\mathop{\mathrm{Hom}}\nolimits}
\def\Ext{\mathop{\mathrm{Ext}}\nolimits}
\def\SheafHom{\mathop{\mathcal{H}\!\mathit{om}}\nolimits}
\def\SheafExt{\mathop{\mathcal{E}\!\mathit{xt}}\nolimits}
\def\Sch{\mathit{Sch}}
\def\Mor{\mathop{\mathrm{Mor}}\nolimits}
\def\Ob{\mathop{\mathrm{Ob}}\nolimits}
\def\Sh{\mathop{\mathit{Sh}}\nolimits}
\def\NL{\mathop{N\!L}\nolimits}
\def\CH{\mathop{\mathrm{CH}}\nolimits}
\def\proetale{{pro\text{-}\acute{e}tale}}
\def\etale{{\acute{e}tale}}
\def\QCoh{\mathit{QCoh}}
\def\Ker{\mathop{\mathrm{Ker}}}
\def\Im{\mathop{\mathrm{Im}}}
\def\Coker{\mathop{\mathrm{Coker}}}
\def\Coim{\mathop{\mathrm{Coim}}}

% Boxtimes
%
\DeclareMathSymbol{\boxtimes}{\mathbin}{AMSa}{"02}

%
% Macros for moduli stacks/spaces
%
\def\QCohstack{\mathcal{QC}\!\mathit{oh}}
\def\Cohstack{\mathcal{C}\!\mathit{oh}}
\def\Spacesstack{\mathcal{S}\!\mathit{paces}}
\def\Quotfunctor{\mathrm{Quot}}
\def\Hilbfunctor{\mathrm{Hilb}}
\def\Curvesstack{\mathcal{C}\!\mathit{urves}}
\def\Polarizedstack{\mathcal{P}\!\mathit{olarized}}
\def\Complexesstack{\mathcal{C}\!\mathit{omplexes}}
% \Pic is the operator that assigns to X its picard group, usage \Pic(X)
% \Picardstack_{X/B} denotes the Picard stack of X over B
% \Picardfunctor_{X/B} denotes the Picard functor of X over B
\def\Pic{\mathop{\mathrm{Pic}}\nolimits}
\def\Picardstack{\mathcal{P}\!\mathit{ic}}
\def\Picardfunctor{\mathrm{Pic}}
\def\Deformationcategory{\mathcal{D}\!\mathit{ef}}
\usepackage{cleveref}
\crefname{originallemma}{original lemma}{original lemmas}
\Crefname{originallemma}{Original lemma}{Original lemmas}
\makeatletter
\newlabel{lemma-good-element}{{00N4}{source}{}{sourcepointer-1}{}}
\newlabel{lemma-good-element@cref}{{[originallemma][0][]00N4}{[1][0][]source}{}{}{}}
\makeatother
\crefname{originallemma}{original lemma}{original lemmas}
\Crefname{originallemma}{Original lemma}{Original lemmas}
\makeatletter
\newlabel{lemma-permute-xi}{{00LJ}{source}{}{sourcepointer-2}{}}
\newlabel{lemma-permute-xi@cref}{{[originallemma][0][]00LJ}{[1][0][]source}{}{}{}}
\makeatother
\crefname{originallemma}{original lemma}{original lemmas}
\Crefname{originallemma}{Original lemma}{Original lemmas}
\makeatletter
\newlabel{lemma-silly}{{00DS}{source}{}{sourcepointer-3}{}}
\newlabel{lemma-silly@cref}{{[originallemma][0][]00DS}{[1][0][]source}{}{}{}}
\makeatother
\title{An element cutting the support down by one acts as a nonzerodivisor on a Cohen-Macaulay module}
\author{The Stacks Project Authors}
\date{stacks-project:algebra.tex:local-snapshot-20261012; GNU FDL 1.2 or later}
\begin{document}
\maketitle
\noindent\textit{Editorial note.} All delivered mathematical text is the author's own, copied from the pinned Stacks Project chapter file algebra.tex; no statement or proof step was rewritten, reordered or removed.\par\medskip
\noindent\textit{Editorial note.} Preamble provenance, as required for this source: the frozen snapshot stages the chapter's own preamble file beside it (preamble.tex, 217 lines, 8116 bytes, the same shared Stacks-wide preamble the published chow and etale-cohomology cards use); the first line of the pinned chapter file is \textbackslash{}input\{preamble\}, so no preamble text lives inside the chapter itself. The card takes that file's lines 8-18 (the verbatim package and the three \textbackslash{}comment-based environments reference/slogan/history at lines 9-12, section 13.2, together with the Xy-pic setup) and its lines 154 to the runtime end of the file (the native theorem declarations and mathematical macros). Not one character was changed, no package was added, no macro was defined, and the excerpt serves only as preamble, never as a span coordinate.\par\medskip
\noindent\textit{Editorial note.} The same-chapter pointers are located at their \textbackslash{}label lines inside the pinned chapter file; the cross-chapter pointers are located in the project's own tag list staged in the snapshot (sources/stacks-chow/tags.txt), at the verbatim lines carrying their canonical tags.\par\medskip
\noindent\textit{Editorial note.} The counted proof refers to three same-chapter statements by label; all are recorded as navigation stubs pointed at their \textbackslash{}label lines, and the arguments consuming them stay inside the counted span.\par\medskip
\noindent\textit{Editorial note.} The counted spans contain no \textbackslash{}cite, so the card carries no bibliography.\par\medskip
\noindent\textit{Allowed standard prerequisites (editorial).}\begin{itemize}
\item Good elements for a maximal regular sequence. Invoked in the counted proof at line 25246. Supplies the case in which the chosen element is a nonzerodivisor with Cohen-Macaulay quotient.
\item Cutting down the dimension by choosing a second element. Used in the induction step of the counted proof (lines 25248-25263). Produces an element controlling the quotient so that the induction hypothesis applies.
\end{itemize}\medskip
\ifdefined\hypertarget\hypertarget{sourcepointer-1}{}\fi
\noindent\textit{Source reference.} Original lemma 00N4: sources/stacks-algebra/pinned-original-source.tex, source lines 25202--25202. Navigation pointer for \textbackslash{}ref\{lemma-good-element\} used inside the counted span (line 25246); the argument that consumes it is retained in the span.\par\medskip
\ifdefined\hypertarget\hypertarget{sourcepointer-2}{}\fi
\noindent\textit{Source reference.} Original lemma 00LJ: sources/stacks-algebra/pinned-original-source.tex, source lines 16856--16856. Navigation pointer for \textbackslash{}ref\{lemma-permute-xi\} used inside the counted span (line 25262); the argument that consumes it is retained in the span.\par\medskip
\ifdefined\hypertarget\hypertarget{sourcepointer-3}{}\fi
\noindent\textit{Source reference.} Original lemma 00DS: sources/stacks-algebra/pinned-original-source.tex, source lines 2781--2781. Navigation pointer for \textbackslash{}ref\{lemma-silly\} used inside the counted span (line 25256); the argument that consumes it is retained in the span.\par\medskip
\begin{lemma}
\label{lemma-CM-one-g}
Let $R$ be a Noetherian local ring.
Let $M$ be a Cohen-Macaulay module over $R$.
Suppose $g \in \mathfrak m$ is such that $\dim(\text{Supp}(M) \cap V(g))
= \dim(\text{Supp}(M)) - 1$. Then (a) $g$ is a nonzerodivisor on $M$,
and (b) $M/gM$ is Cohen-Macaulay of depth one less.
\end{lemma}

\begin{proof}
Choose a $M$-regular sequence $f_1, \ldots, f_d$ with
$d = \dim(\text{Supp}(M))$. If $g$ is good with respect to
$(M, f_1, \ldots, f_d)$ we win by Lemma \ref{lemma-good-element}.
In particular the lemma holds if $d = 1$. (The case $d = 0$ does
not occur.) Assume $d > 1$. Choose an element $h \in R$ such that
(\romannumeral1) $h$ is good with respect to $(M, f_1, \ldots, f_d)$,
and (\romannumeral2) $\dim(\text{Supp}(M) \cap V(h, g)) = d - 2$.
To see $h$ exists, let $\{\mathfrak q_j\}$ be the (finite) set of
minimal primes of the closed sets $\text{Supp}(M)$,
$\text{Supp}(M)\cap V(f_1, \ldots, f_i)$, $i = 1, \ldots, d - 1$,
and $\text{Supp}(M) \cap V(g)$. None of these $\mathfrak q_j$
is equal to $\mathfrak m$ and hence we may find $h \in \mathfrak m$,
$h \not \in \mathfrak q_j$ by Lemma \ref{lemma-silly}. It is clear
that $h$ satisfies (\romannumeral1) and (\romannumeral2). From
Lemma \ref{lemma-good-element} we conclude that
$M/hM$ is Cohen-Macaulay. By (\romannumeral2) we see that the pair
$(M/hM, g)$ satisfies the induction hypothesis. Hence
$M/(h, g)M$ is Cohen-Macaulay and $g : M/hM \to M/hM$
is injective. By Lemma \ref{lemma-permute-xi} we see
that $g : M \to M$ and $h : M/gM \to M/gM$
are injective. Combined with the fact that $M/(g, h)M$
is Cohen-Macaulay this finishes the proof.
\end{proof}

\end{document}
